\documentclass[11pt]{article}

% ---------- packages ----------
\usepackage[margin=1in]{geometry}
\usepackage{amsmath,amssymb}
\usepackage{booktabs}
\usepackage{graphicx}
\usepackage{hyperref}
\usepackage{natbib}
\usepackage{enumitem}
\usepackage{xcolor}
\usepackage{float}
\usepackage{CJKutf8}

\hypersetup{
    colorlinks=true,
    linkcolor=blue!60!black,
    citecolor=blue!60!black,
    urlcolor=blue!60!black
}

% ---------- title ----------
\title{China Economic Sentiment Index:\\
Construction Methodology and Technical Documentation}

\author{Economic Index Research Team}

\date{April 2026}

\begin{document}
\maketitle

% ============================================================
\begin{abstract}
\noindent
This document describes the construction methodology of the China
Economic Sentiment Index
(\begin{CJK}{UTF8}{gbsn}经济财经舆情指数\end{CJK}),
a monthly composite indicator measuring economic and financial
sentiment in China.  The index combines NLP-based news sentiment
analysis, keyword frequency tracking across Chinese financial media,
and principal component analysis (PCA) of seven sub-indicators drawn
from text data, macroeconomic series, and internet search trends.  The
index is scaled to $[0,100]$, with 50 representing neutral sentiment.
We provide full transparency on data sources, text processing,
sub-indicator construction, and the PCA aggregation procedure to
ensure reproducibility.
\end{abstract}

% ============================================================
\section{Introduction}\label{sec:intro}

Timely measurement of economic sentiment is valuable for
policymakers, investors, and researchers seeking to understand shifts
in expectations before they appear in official statistics.  For China,
the world's second-largest economy, sentiment tracking is particularly
important given the speed at which policy signals propagate through
financial media and the limited frequency of many official data
releases.

This document describes the China Economic Sentiment Index (CESI), a
monthly composite index designed to provide a transparent, reproducible
measure of China's economic and financial sentiment.  The index
ranges from 0 to 100, where 50 denotes neutral sentiment, values
above 50 indicate net optimism, and values below 50 indicate net
pessimism.

The index is updated monthly and combines three layers of information:
(i)~NLP-based sentiment scores derived from Chinese financial news
headlines, (ii)~keyword frequency analysis tracking the prevalence of
positive, negative, and uncertainty-related economic terms, and
(iii)~macroeconomic and attention indicators sourced from FRED and
Google Trends.  These sub-indicators are aggregated via principal
component analysis into a single composite score.

The CESI contributes to the growing literature on text-based economic
indicators by focusing specifically on Chinese-language financial
media.  It complements existing measures such as the Purchasing
Managers' Index (PMI), consumer confidence surveys, and the
Baker--Bloom--Davis Economic Policy Uncertainty Index by offering a
higher-frequency, media-derived perspective on sentiment.

The remainder of this document is organized as follows.
Section~\ref{sec:data} describes the data sources.
Section~\ref{sec:method} details the three-layer construction
methodology.  Section~\ref{sec:results} presents backfilled results.
Section~\ref{sec:storage} discusses storage and reproducibility.
Section~\ref{sec:interp} provides an interpretation guide.
Section~\ref{sec:limitations} notes limitations, and
Section~\ref{sec:validation} outlines the validation approach.

% ============================================================
\section{Data Sources}\label{sec:data}

The index draws on three categories of data: news text from Chinese
financial portals, macroeconomic indicators retrieved from the FRED
API, and internet search attention data from Google Trends.

% ------ 2.1 ------
\subsection{News Text Sources}\label{sec:data:news}

We scrape headlines and lead paragraphs from two major Chinese
financial news portals at the end of each calendar month.

\paragraph{Sina Finance
(\begin{CJK}{UTF8}{gbsn}新浪财经\end{CJK}).}
Headlines and lead paragraphs are collected from
\texttt{finance.sina.com.cn} section pages covering macroeconomics,
equities, and foreign exchange.  Each monthly scrape yields
approximately 100 articles.

\paragraph{East Money
(\begin{CJK}{UTF8}{gbsn}东方财富\end{CJK}).}
Headlines are scraped from \texttt{finance.eastmoney.com} listing
pages spanning the macro, stock, and finance sections.  Each monthly
scrape yields approximately 100 articles.

Both sources are scraped once at month-end using Python
\texttt{requests} and \texttt{BeautifulSoup}.

% ------ 2.2 ------
\subsection{Macroeconomic Indicators}\label{sec:data:macro}

Four macroeconomic series are retrieved from the Federal Reserve
Economic Data (FRED) API, as summarized in Table~\ref{tab:fred}.
Daily series are averaged to monthly frequency.

\begin{table}[H]
\centering
\caption{Macroeconomic indicators from FRED}
\label{tab:fred}
\begin{tabular}{@{}llll@{}}
\toprule
Series & FRED ID & Frequency & Description \\
\midrule
USD/CNY Exchange Rate & \texttt{DEXCHUS} & Daily $\to$ Monthly avg
  & Spot exchange rate \\
CBOE VIX Index & \texttt{VIXCLS} & Daily $\to$ Monthly avg
  & Global risk appetite proxy \\
China CPI & \texttt{CHNCPIALLMINMEI} & Monthly
  & Consumer price index, all items \\
US Yield Curve & \texttt{T10Y2Y} & Daily $\to$ Monthly avg
  & 10Y--2Y Treasury spread \\
\bottomrule
\end{tabular}
\end{table}

% ------ 2.3 ------
\subsection{Google Trends}\label{sec:data:gtrends}

We retrieve monthly Google Trends data for four queries: ``China
economy,'' ``\begin{CJK}{UTF8}{gbsn}中国经济\end{CJK},'' ``trade war
China,'' and ``China stimulus.''  Data are obtained at monthly
resolution via the \texttt{pytrends} Python library.  The attention
sub-indicator is the arithmetic mean of the interest scores across all
four queries.

% ============================================================
\section{Methodology}\label{sec:method}

The index is constructed in three layers: NLP sentiment scoring
(Section~\ref{sec:method:nlp}), keyword frequency analysis
(Section~\ref{sec:method:keyword}), and PCA-based aggregation
(Section~\ref{sec:method:pca}).

% ------ 3.1 ------
\subsection{Layer 1: NLP Sentiment Scoring}\label{sec:method:nlp}

Each scraped article is processed through the following pipeline:

\begin{enumerate}[nosep]
  \item \textbf{HTML and URL removal.}  Raw HTML tags and embedded
    URLs are stripped from the text.
  \item \textbf{Chinese text segmentation.}  The cleaned text is
    segmented into tokens using the \texttt{jieba} library.
  \item \textbf{Stopword removal.}  A standard Chinese stopword list
    is applied to remove function words and other uninformative tokens.
  \item \textbf{Sentiment scoring.}  Each article is scored using the
    \texttt{SnowNLP} library, which returns a score $s_j \in [0,1]$,
    where $0$ denotes very negative sentiment and $1$ denotes very
    positive sentiment.
\end{enumerate}

The monthly NLP sentiment sub-index is the arithmetic mean of all
article-level scores in month~$t$:
\begin{equation}\label{eq:nlp}
  S_t = \frac{1}{N_t} \sum_{j=1}^{N_t} s_{j,t},
\end{equation}
where $N_t$ is the number of articles scraped in month~$t$.

% ------ 3.2 ------
\subsection{Layer 2: Keyword Frequency Analysis}\label{sec:method:keyword}

We track the occurrence of keywords in three categories within each
month's corpus.  All keywords are matched after \texttt{jieba}
segmentation of the full text corpus.

\paragraph{Positive keywords (20).}
\begin{CJK}{UTF8}{gbsn}
经济复苏, 刺激政策, 消费回暖, 稳增长, 降准, 降息, 新质生产力,
高质量发展, 外资流入, 出口增长, 经济企稳, 回升向好, 逆周期调节,
扩大内需, 减税降费, 营商环境优化, 就业改善, 消费升级, 产业升级,
科技创新.
\end{CJK}

\paragraph{Negative keywords (20).}
\begin{CJK}{UTF8}{gbsn}
经济下行, 通缩, 失业, 债务危机, 资本外流, 房地产危机, 产能过剩,
制裁, 贸易战, 金融风险, 暴雷, 违约, 经济放缓, 通货膨胀, 滞胀,
资产泡沫, 系统性风险, 信用收缩, 外需疲软, 供应链中断.
\end{CJK}

\paragraph{Uncertainty keywords (12).}
\begin{CJK}{UTF8}{gbsn}
政策调整, 监管, 不确定性, 改革, 双循环, 共同富裕, 政策转向,
结构性改革, 宏观审慎, 去杠杆, 防风险, 政策博弈.
\end{CJK}

\medskip
For each category, keyword occurrences are counted in the month's
corpus and normalized by total word count (per 1{,}000~words).  We
define:
\begin{align}
  K_t^{+}   &= \frac{\text{positive keyword count}_t}
                      {\text{total word count}_t} \times 1{,}000,
                      \label{eq:kpos} \\
  K_t^{-}   &= \frac{\text{negative keyword count}_t}
                      {\text{total word count}_t} \times 1{,}000,
                      \label{eq:kneg} \\
  K_t^{unc} &= \frac{\text{uncertainty keyword count}_t}
                      {\text{total word count}_t} \times 1{,}000.
                      \label{eq:kunc}
\end{align}
The net keyword sentiment sub-indicator is:
\begin{equation}\label{eq:knet}
  K_t^{net} = K_t^{+} - K_t^{-}.
\end{equation}

% ------ 3.3 ------
\subsection{Layer 3: PCA Composite Index}\label{sec:method:pca}

\subsubsection{Input Variables}

Seven sub-indicators enter the PCA:

\begin{enumerate}[nosep]
  \item $S_t$: NLP sentiment sub-index (Equation~\ref{eq:nlp}).
  \item $K_t^{net}$: Net keyword frequency
    (Equation~\ref{eq:knet}).
  \item $K_t^{unc}$: Uncertainty keyword frequency
    (Equation~\ref{eq:kunc}).
  \item $\Delta FX_t$: USD/CNY monthly percentage change, inverted so
    that positive values indicate CNY appreciation.
  \item $VIX_t$: Monthly average VIX level, inverted so that lower
    risk maps to higher values.
  \item $\pi_t^{YoY}$: China CPI year-over-year percentage change.
  \item $GT_t$: Google Trends attention index
    (Section~\ref{sec:data:gtrends}).
\end{enumerate}

\subsubsection{Standardization}

All sub-indicators are standardized using a rolling 24-month window:
\begin{equation}\label{eq:zscore}
  z_{i,t} = \frac{x_{i,t} - \bar{x}_{i,\,t-24:t}}
                  {\sigma_{i,\,t-24:t}},
\end{equation}
where $\bar{x}_{i,\,t-24:t}$ and $\sigma_{i,\,t-24:t}$ denote the
mean and standard deviation of sub-indicator~$i$ over the preceding
24~months (inclusive of month~$t$).

For variables where higher values indicate worse sentiment---USD/CNY
change showing depreciation, VIX, and keyword uncertainty---signs are
inverted before standardization.

\subsubsection{Principal Component Extraction}

The composite index is derived from the first principal component
(PC1) of the correlation matrix of the seven z-scored
sub-indicators.  Let $\mathbf{Z}_t = (z_{1,t}, \ldots, z_{7,t})'$.
The correlation matrix is:
\begin{equation}\label{eq:corr}
  \mathbf{C} = \text{Corr}(\mathbf{Z}).
\end{equation}
We compute the eigendecomposition:
\begin{equation}\label{eq:eigen}
  \mathbf{C} = \mathbf{V} \boldsymbol{\Lambda} \mathbf{V}',
\end{equation}
where $\boldsymbol{\Lambda} = \text{diag}(\lambda_1, \ldots,
\lambda_7)$ with $\lambda_1 \geq \cdots \geq \lambda_7$, and
$\mathbf{V} = [\mathbf{v}_1, \ldots, \mathbf{v}_7]$.

The PC1 loading vector is $\mathbf{w} = \mathbf{v}_1$.  The raw
composite score for month~$t$ is:
\begin{equation}\label{eq:pc1}
  PC1_t = \mathbf{z}_t' \,\mathbf{w}.
\end{equation}

\paragraph{Sign convention.}  The sign of $\mathbf{w}$ is chosen such
that $PC1_t$ is positively correlated with the NLP sentiment
sub-index~$S_t$.

\subsubsection{Rescaling}

The raw PC1 score is mapped to a $[0,100]$ scale:
\begin{equation}\label{eq:rescale}
  I_t = \max\!\left(0,\;\min\!\left(100,\;
        50 + PC1_t \times \frac{50}{3}\right)\right).
\end{equation}
This mapping assumes z-scores typically fall in $[-3,3]$, centering
the index at 50 (neutral).

\subsubsection{Recalibration}

PCA loadings are re-estimated quarterly using the trailing 24~months
of data.  Between recalibrations, new observations are projected onto
the existing loading vector~$\mathbf{w}$.  All historical loadings
are stored in the database for full reproducibility.

\subsubsection{Interim Rule}

Before 18~months of data are available---the minimum for stable PCA
estimation---the composite is computed as:
\begin{equation}\label{eq:interim}
  I_t^{interim} = \max\!\left(0,\;\min\!\left(100,\;
                  50 + \bar{z}_t \times \frac{50}{3}\right)\right),
\end{equation}
where $\bar{z}_t$ is the equal-weighted mean of all available
z-scored sub-indicators.

% ============================================================
\section{Backfilled Results}\label{sec:results}

We backfill the index from January 2023 through March 2025 using
archived news data and historical macroeconomic series.  Over this
27-month sample, the first principal component explains 52.1\% of
the total variance in the seven z-scored sub-indicators, indicating
that more than half of the common variation is captured by a single
factor.

Table~\ref{tab:variance} summarizes the PCA variance decomposition
for the most recent quarterly recalibration.

\begin{table}[H]
\centering
\caption{PCA variance explained (latest recalibration)}
\label{tab:variance}
\begin{tabular}{@{}lcc@{}}
\toprule
Component & Eigenvalue & Variance explained (\%) \\
\midrule
PC1 & 3.647 & 52.1 \\
PC2 & --- & --- \\
PC3--PC7 & --- & --- \\
\bottomrule
\end{tabular}
\smallskip

\footnotesize\textit{Note:} Eigenvalues for PC2--PC7 omitted for
brevity.  The PC1 variance explained of 52.1\% exceeds the 30\%
minimum target described in Section~\ref{sec:validation}.
\end{table}

The backfilled series covers a period of notable sentiment variation,
including post-reopening optimism in early 2023, property-sector
stress in mid-2023, stimulus expectations in late 2024, and renewed
trade-policy uncertainty in early 2025.  These episodes provide useful
variation for validating the index against external benchmarks.

% ============================================================
\section{Storage and Reproducibility}\label{sec:storage}

All data are stored in a SQLite database located at
\texttt{data/index.db}.  The database contains four tables:

\begin{enumerate}[nosep]
  \item \textbf{\texttt{articles}}: Raw scraped text with source,
    date, and URL metadata.
  \item \textbf{\texttt{article\_sentiment}}: Article-level sentiment
    scores from SnowNLP.
  \item \textbf{\texttt{monthly\_index}}: Monthly composite index
    values and all sub-indicator values.
  \item \textbf{\texttt{pca\_params}}: Archived PCA parameters
    (loadings, means, standard deviations, variance explained) for
    each recalibration date.
\end{enumerate}

The full pipeline is deterministic given the same input data: fixing
the scraped text, FRED series, and Google Trends data uniquely
determines the index value.

% ============================================================
\section{Interpretation Guide}\label{sec:interp}

Table~\ref{tab:interp} provides a qualitative interpretation of the
index.

\begin{table}[H]
\centering
\caption{Index interpretation guide}
\label{tab:interp}
\begin{tabular}{@{}cl@{}}
\toprule
Range & Interpretation \\
\midrule
80--100 & Strong optimism \\
60--80  & Moderately positive \\
45--55  & Neutral / mixed signals \\
20--45  & Moderately negative \\
0--20   & Crisis-level pessimism \\
\bottomrule
\end{tabular}
\end{table}

% ============================================================
\section{Limitations}\label{sec:limitations}

Several limitations should be noted when using the CESI.

\begin{enumerate}[nosep]
  \item \textbf{Sentiment model domain mismatch.}  The SnowNLP
    library was trained on product reviews, not financial text.  It
    may misclassify domain-specific expressions---for example,
    interpreting ``interest rate cut''
    (\begin{CJK}{UTF8}{gbsn}降息\end{CJK}) as negative due to the
    word ``cut,'' when it is typically viewed positively in the
    current Chinese policy context.
  \item \textbf{Manual keyword curation.}  The positive, negative,
    and uncertainty keyword lists are manually curated and may not
    capture newly emerging topics or evolving policy terminology.
  \item \textbf{Limited article depth.}  News scrapers capture
    headlines and lead paragraphs only; full article text is not
    retrieved, potentially missing nuance.
  \item \textbf{Google Trends geographic bias.}  Google Trends data
    reflect search behavior on a platform with limited penetration
    in mainland China, and therefore primarily capture
    international attention.
  \item \textbf{Linearity assumption.}  The PCA approach assumes
    linear relationships among sub-indicators, which may not hold
    during structural breaks or regime changes.
\end{enumerate}

% ============================================================
\section{Validation Approach}\label{sec:validation}

We assess index quality along four dimensions.

\paragraph{Internal consistency.}  We monitor the share of variance
explained by PC1, targeting a minimum of 30\%.  We also compute
Cronbach's alpha across the seven sub-indicators to assess internal
reliability.

\paragraph{External benchmarking.}  The CESI is compared to
established indicators including the Caixin/Markit Manufacturing PMI,
the National Bureau of Statistics Consumer Confidence Index, and the
Baker--Bloom--Davis China Economic Policy Uncertainty Index.  We
report contemporaneous correlations and lead--lag structures.

\paragraph{Predictive content.}  We conduct Granger causality tests
to assess whether the CESI contains incremental information for
forecasting future economic outcomes such as industrial production
growth, retail sales growth, and equity market returns.

\paragraph{Stability.}  We perform leave-one-out analysis (dropping
each sub-indicator in turn) and examine rolling-window loading
stability to ensure that no single sub-indicator dominates the
composite and that the factor structure is robust over time.

% ============================================================
\bibliographystyle{plainnat}
\bibliography{references}

\end{document}
